\chapter{Chapter04 — Circuit Influence}
%%% generated by the blueprint pipeline from FABL.Chapter04.CircuitInfluence
\section{Overview}

Declarations of the Lean module \texttt{FABL.Chapter04.CircuitInfluence}.

\section{Declarations}

\begin{definition}[Width bound for a gate]
\lean{FABL.CircuitGate.HasWidthLE}
\statementsource{Analysis-of-Boolean-Functions-by-Ryan-ODonnell}{
  pdf-ffbe9c7b3306-p093-b004,
  pdf-ffbe9c7b3306-p094-b002,
  pdf-ffbe9c7b3306-p103-b002,
  pdf-ffbe9c7b3306-p104-b001
}
\label{FABL.CircuitGate.HasWidthLE}
\leanok
\uses{FABL.BooleanFunction, FABL.CircuitGate, FABL.HasCNFWidthLE, FABL.HasDNFWidthLE}
Let $f$ be a Boolean function on $n$ sign variables, let $w\in\mathbb{N}$, and let the
gate be either AND or OR. The gate has width at most $w$ for $f$ if, when it is AND, $f$
has a conjunctive normal form of width at most $w$, and, when it is OR, $f$ has a
disjunctive normal form of width at most $w$.
\end{definition}

\begin{theorem}[Monotonicity of representation width]
\lean{FABL.CircuitGate.hasWidthLE_mono}
\label{FABL.CircuitGate.hasWidthLE_mono}
\leanok
\uses{FABL.BooleanFunction, FABL.CircuitGate, FABL.CircuitGate.HasWidthLE}
\difficulty{4}
Let $n$ be a nonnegative integer, let $f:\{\pm1\}^n\to\{\pm1\}$, let $G$ be an AND or OR
gate, and let $u,v$ be nonnegative integers with $u\le v$. If $f$ has width at most $u$
in the normal form associated with $G$, then it has width at most $v$ in that normal
form.
\end{theorem}

\begin{definition}[Pointwise evaluation of a Boolean gate]
\lean{FABL.CircuitGate.evalFinsetFunction}
\label{FABL.CircuitGate.evalFinsetFunction}
\leanok
\uses{FABL.BooleanFunction, FABL.CircuitGate, FABL.CircuitGate.evalFinset}
Given a logical AND or OR gate, a finite family of Boolean functions, and a finite set
of input indices, the resulting Boolean function applies the gate to the selected
functions’ values at each input point.
\end{definition}

\begin{theorem}[Evaluation of a gate on functions]
\lean{FABL.CircuitGate.evalFinsetFunction_apply}
\label{FABL.CircuitGate.evalFinsetFunction_apply}
\leanok
\uses{FABL.BooleanFunction, FABL.CircuitGate, FABL.CircuitGate.evalFinset, FABL.CircuitGate.evalFinsetFunction}
\difficulty{1}
Let $m,n$ be natural numbers, let a logical AND or OR gate be given, and let $m$ Boolean
functions on the $n$-dimensional sign cube take values in $\\{ -1,1 \\}$. For any subset
of the $m$ input indices, evaluating the gate pointwise at a cube point $x$ gives the
same result as applying the gate to the values of those functions at $x$.
\end{theorem}

\begin{theorem}[Closure under a finite gate operation]
\lean{FABL.CircuitGate.hasWidthLE_evalFinsetFunction}
\label{FABL.CircuitGate.hasWidthLE_evalFinsetFunction}
\leanok
\uses{FABL.BooleanFunction, FABL.CNFFormula, FABL.CNFFormula.clauseEval, FABL.CNFFormula.eval, FABL.CNFFormula.toBooleanFunction, FABL.CNFFormula.width, FABL.CircuitGate, FABL.CircuitGate.HasWidthLE, FABL.CircuitGate.evalFinset, FABL.CircuitGate.evalFinsetFunction, FABL.CircuitGate.evalFinset_and_eq_neg_one_iff, FABL.CircuitGate.evalFinset_or_eq_neg_one_iff, FABL.DNFFormula, FABL.DNFFormula.eval, FABL.DNFFormula.eval_eq_neg_one_iff, FABL.DNFFormula.toBooleanFunction, FABL.DNFFormula.width, FABL.DNFFormula.width_le_iff, FABL.DNFFormula.width_le_of_mem}
\difficulty{5}
Let $m,n$ be nonnegative integers, let $G$ be an AND or OR gate, let
$v_i:\{\pm1\}^n\to\{\pm1\}$ be a family of functions indexed by $i\in\text{Fin}(m)$, and
let $S\subseteq\text{Fin}(m)$. Suppose each $v_i$ with $i\in S$ has width at most $w$ in
the normal form associated with $G$. The function obtained by applying $G$ pointwise to
the functions indexed by $S$ also has width at most $w$ in that normal form.
\end{theorem}

\begin{theorem}[DNF width from decision tree depth]
\lean{FABL.F₂DecisionTree.hasDNFWidthLE_decisionTreeDepth}
\leanok
\uses{FABL.BooleanFunction, FABL.F₂DecisionTree.decisionTreeDepth, FABL.F₂DecisionTree.exists_computingTree_depth_eq_decisionTreeDepth, FABL.F₂DecisionTree.hasDNFSizeWidth_of_computes, FABL.HasDNFWidthLE, FABL.binaryOfBooleanFunction}
\proofsource{Analysis-of-Boolean-Functions-by-Ryan-ODonnell}{pdf-ffbe9c7b3306-p094-b006}
\difficulty{4}
Let $n$ be a nonnegative integer and let $f:\{−1,1\}^n\to\{−1,1\}$ be a Boolean
function. If $d$ is the least depth of a decision tree on $\mathbb{F}_2^n$ computing the
coordinatewise binary encoding of $f$, then $f$ has a disjunctive normal form in which
every term contains at most $d$ literals.
\end{theorem}

\begin{theorem}[Injectivity of Boolean duality]
\lean{FABL.booleanDual_injective}
\label{FABL.booleanDual_injective}
\leanok
\uses{FABL.BooleanFunction, FABL.CNFFormula.booleanDual}
\difficulty{3}
Let $n$ be a nonnegative integer. For Boolean functions
$f,g\colon\{\pm1\}^n\to\{\pm1\}$, if $-f(-x)=-g(-x)$ for every $x\in\{\pm1\}^n$, then
$f(x)=g(x)$ for every such $x$, where $-x$ denotes coordinatewise negation.
\end{theorem}

\begin{theorem}[Has cnf width le decision tree depth]
\lean{FABL.F₂DecisionTree.hasCNFWidthLE_decisionTreeDepth}
\leanok
\uses{FABL.BooleanFunction, FABL.CNFFormula, FABL.CNFFormula.booleanDual, FABL.CNFFormula.switchAndOr, FABL.CNFFormula.switchAndOr_toBooleanFunction, FABL.CNFFormula.toBooleanFunction, FABL.CNFFormula.width, FABL.DNFFormula.toBooleanFunction, FABL.F₂DecisionTree.decisionTreeDepth, FABL.F₂DecisionTree.hasDNFWidthLE_decisionTreeDepth, FABL.HasCNFWidthLE, FABL.binaryBooleanDual, FABL.binaryOfBooleanFunction, FABL.booleanDual_injective, FABL.decisionTreeDepth_binaryBooleanDual, FABL.signEncode_binaryCubeComplement}
\statementsource{Analysis-of-Boolean-Functions-by-Ryan-ODonnell}{pdf-ffbe9c7b3306-p094-b006}
\difficulty{5}
% dataset-pipeline: stub -- prose pending build_informal_statements.py
Statement pending regeneration.
\end{theorem}

\begin{theorem}[Normal-form width bounded by decision-tree depth]
\lean{FABL.CircuitGate.hasWidthLE_decisionTreeDepth}
\proofsource{Analysis-of-Boolean-Functions-by-Ryan-ODonnell}{pdf-ffbe9c7b3306-p094-b006}
\label{FABL.CircuitGate.hasWidthLE_decisionTreeDepth}
\leanok
\uses{FABL.BooleanFunction, FABL.CircuitGate, FABL.CircuitGate.HasWidthLE, FABL.F₂DecisionTree.decisionTreeDepth, FABL.F₂DecisionTree.hasCNFWidthLE_decisionTreeDepth, FABL.F₂DecisionTree.hasDNFWidthLE_decisionTreeDepth, FABL.binaryOfBooleanFunction}
\difficulty{4}
Let $n$ be a nonnegative integer, let $f:\{\pm1\}^n\to\{\pm1\}$, and let $G$ be an AND
or OR gate. Then $f$ has a representation of width at most the decision-tree depth of
the function on binary inputs obtained by converting each binary coordinate to its
corresponding sign.
\end{theorem}

\begin{theorem}[Influence bound for bounded-width formulas]
\lean{FABL.totalInfluence_le_two_mul_of_hasWidthLE}
\proofsource{Analysis-of-Boolean-Functions-by-Ryan-ODonnell}{
  pdf-ffbe9c7b3306-p095-b002,
  pdf-ffbe9c7b3306-p095-b003
}
\label{FABL.totalInfluence_le_two_mul_of_hasWidthLE}
\leanok
\uses{FABL.BooleanFunction, FABL.BooleanFunction.toReal, FABL.CircuitGate, FABL.CircuitGate.HasWidthLE, FABL.totalInfluence, FABL.totalInfluence_le_two_mul_of_hasCNFWidthLE, FABL.totalInfluence_le_two_mul_of_hasDNFWidthLE}
\difficulty{4}
Let $n$ be a natural number, let $f$ be a Boolean function on the $n$-dimensional sign
cube, and choose either the AND gate or the OR gate. If the chosen gate has a formula
representation of width at most $w$—a CNF representation for AND or a DNF representation
for OR—then the total influence of $f$ is at most $2w$.
\end{theorem}

\begin{theorem}[Width-zero gate functions are constant]
\lean{FABL.CircuitGate.exists_eq_const_of_hasWidthLE_zero}
\label{FABL.CircuitGate.exists_eq_const_of_hasWidthLE_zero}
\leanok
\uses{FABL.BooleanFunction, FABL.CNFFormula.booleanDual, FABL.CNFFormula.switchAndOr, FABL.CNFFormula.switchAndOr_toBooleanFunction, FABL.CNFFormula.width, FABL.CircuitGate, FABL.CircuitGate.HasWidthLE, FABL.DNFFormula.toBooleanFunction_eq_of_width_eq_zero, FABL.DNFFormula.width, FABL.Sign}
\difficulty{5}
Let $n$ be a nonnegative integer, let $f:\{\pm1\}^n\to\{\pm1\}$, and let $G$ be an AND
or OR gate. Suppose $f$ has a representation of width at most zero in the normal form
associated with $G$ (CNF for AND and DNF for OR). Then $f$ is constant: there is a sign
$c\in\{\pm1\}$ such that $f(x)=c$ for every input $x$.
\end{theorem}

\begin{definition}[Evaluation of a circuit tail]
\lean{FABL.CircuitTail.evalFunction}
\label{FABL.CircuitTail.evalFunction}
\leanok
\uses{FABL.BooleanFunction, FABL.CircuitGate, FABL.CircuitTail, FABL.CircuitTail.eval}
Given a circuit tail, a logical AND or OR gate, and a family of Boolean functions, the
resulting Boolean function evaluates the tail at each input using the supplied functions
as its initial values.
\end{definition}

\begin{theorem}[Evaluation of an output layer]
\lean{FABL.CircuitTail.evalFunction_output}
\label{FABL.CircuitTail.evalFunction_output}
\leanok
\uses{FABL.BooleanFunction, FABL.CircuitGate, FABL.CircuitGate.evalFinsetFunction, FABL.CircuitTail.evalFunction}
\difficulty{1}
Let $m,n$ be natural numbers, let an output layer be specified by a finite set of input
indices among $m$ indices, choose an AND or OR gate, and assign a Boolean function on
the $n$-dimensional sign cube to each input. Evaluating this output layer as a function
gives the result of applying the chosen gate to those input functions at the selected
indices.
\end{theorem}

\begin{theorem}[Evaluation of a layered circuit tail]
\lean{FABL.CircuitTail.evalFunction_layer}
\label{FABL.CircuitTail.evalFunction_layer}
\leanok
\uses{FABL.BooleanFunction, FABL.CircuitGate, FABL.CircuitGate.dual, FABL.CircuitGate.evalFinsetFunction, FABL.CircuitTail, FABL.CircuitTail.evalFunction}
\difficulty{1}
Let $m,n$ be natural numbers, let a circuit tail consist of a list of gates followed by
a remaining tail, choose an AND or OR gate, and assign a Boolean function on the
$n$-dimensional sign cube to each of the $m$ inputs. Evaluating the whole tail as a
function is the same as evaluating the remaining tail with the dual gate on the
functions obtained by applying the chosen gate to each listed input set.
\end{theorem}

\begin{theorem}[Extended sign restriction eval function]
\lean{FABL.CircuitTail.extendedSignRestriction_evalFunction}
\statementsource{Analysis-of-Boolean-Functions-by-Ryan-ODonnell}{
  pdf-ffbe9c7b3306-p098-b003,
  pdf-ffbe9c7b3306-p098-b004
}
\label{FABL.CircuitTail.extendedSignRestriction_evalFunction}
\leanok
\uses{FABL.BooleanFunction, FABL.CircuitGate, FABL.CircuitTail, FABL.CircuitTail.evalFunction, FABL.FixedSignCube, FABL.extendedSignRestriction}
\difficulty{3}
% dataset-pipeline: stub -- prose pending build_informal_statements.py
Statement pending regeneration.
\end{theorem}

\begin{theorem}[Constant inputs yield constant output]
\lean{FABL.CircuitTail.exists_eq_const_evalFunction_of_values}
\label{FABL.CircuitTail.exists_eq_const_evalFunction_of_values}
\leanok
\uses{FABL.BooleanFunction, FABL.CircuitGate, FABL.CircuitGate.dual, FABL.CircuitGate.evalFinset, FABL.CircuitGate.evalFinsetFunction, FABL.CircuitTail, FABL.CircuitTail.eval, FABL.CircuitTail.evalFunction, FABL.CircuitTail.evalFunction_layer, FABL.Sign}
\difficulty{4}
Let $T$ be a circuit tail with $m$ inputs, let $g$ be an AND or OR gate, and let each
input be assigned a Boolean function on $-1,1^n$ taking values in $\text{Sign}$.
Suppose each such function is constant, with its own constant sign. Then evaluating the
tail pointwise gives a constant Boolean function on $-1,1^n$.
\end{theorem}

\begin{theorem}[Zero influence for width-zero inputs]
\lean{FABL.CircuitTail.totalInfluence_evalFunction_eq_zero_of_width_zero}
\label{FABL.CircuitTail.totalInfluence_evalFunction_eq_zero_of_width_zero}
\leanok
\uses{FABL.BooleanFunction, FABL.BooleanFunction.toReal, FABL.CircuitGate, FABL.CircuitGate.HasWidthLE, FABL.CircuitGate.exists_eq_const_of_hasWidthLE_zero, FABL.CircuitTail, FABL.CircuitTail.evalFunction, FABL.CircuitTail.exists_eq_const_evalFunction_of_values, FABL.Sign, FABL.signValue, FABL.totalInfluence, FABL.totalInfluence_const}
\difficulty{4}
Let $T$ be a circuit tail with $m$ inputs, let $g$ be an AND or OR gate, and assign a
sign-valued Boolean function to each input on the $n$-dimensional sign cube. Suppose
every input function has a gate-appropriate conjunctive or disjunctive normal form of
width at most zero. Then the total influence of the real-valued function obtained by
evaluating $T$ is zero.
\end{theorem}

\begin{definition}[Maximum restricted decision tree depth]
\lean{FABL.maxRestrictedDecisionTreeDepth}
\label{FABL.maxRestrictedDecisionTreeDepth}
\leanok
\uses{FABL.BooleanFunction, FABL.FixedSignCube, FABL.restrictedDecisionTreeDepth}
Given a finite family of Boolean functions, a set $J$ of coordinates, and a sign
assignment $z$ on the coordinates outside $J$, the maximum restricted decision tree
depth is the maximum, over the family, of the functions’ restricted decision tree depths
with $J$ and $z$ fixed.
\end{definition}

\begin{theorem}[Restricted decision tree depth bound]
\lean{FABL.restrictedDecisionTreeDepth_le_max}
\label{FABL.restrictedDecisionTreeDepth_le_max}
\leanok
\uses{FABL.BooleanFunction, FABL.FixedSignCube, FABL.maxRestrictedDecisionTreeDepth, FABL.restrictedDecisionTreeDepth}
\difficulty{1}
Let $n,m\in\mathbb N$, let $f_0,\ldots,f_{m-1}$ be Boolean functions on the sign cube of
dimension $n$, let $J$ be a set of coordinates, and fix signs on the coordinates outside
$J$. For each $i<m$, the decision tree depth of the restriction of $f_i$ with those
fixed signs is at most the maximum of these restricted depths over all $i<m$.
\end{theorem}

\begin{theorem}[Restricted functions have bounded gate width]
\lean{FABL.CircuitGate.hasWidthLE_extendedSignRestriction_max}
\label{FABL.CircuitGate.hasWidthLE_extendedSignRestriction_max}
\leanok
\uses{FABL.BooleanFunction, FABL.CircuitGate, FABL.CircuitGate.HasWidthLE, FABL.CircuitGate.hasWidthLE_decisionTreeDepth, FABL.CircuitGate.hasWidthLE_mono, FABL.FixedSignCube, FABL.F₂DecisionTree.decisionTreeDepth, FABL.binaryOfBooleanFunction, FABL.extendedSignRestriction, FABL.maxRestrictedDecisionTreeDepth, FABL.restrictedDecisionTreeDepth, FABL.restrictedDecisionTreeDepth_le_max}
\difficulty{3}
Let $m,n$ be nonnegative integers, let $G$ be an AND or OR gate, and let
$v_i:\{\pm1\}^n\to\{\pm1\}$ be a family indexed by $i\in\text{Fin}(m)$. Choose a
coordinate set $J\subseteq\text{Fin}(n)$, an assignment of signs to the coordinates
outside $J$, and an index $i\in\text{Fin}(m)$. The restriction of $v_i$ obtained by
fixing those outside coordinates has width, in the normal form associated with $G$, at
most the maximum decision-tree depth of the similarly restricted functions in the
family.
\end{theorem}

\begin{theorem}[Exponential bound from maximum depth]
\lean{FABL.two_pow_sup_le_one_add_sum}
\label{FABL.two_pow_sup_le_one_add_sum}
\leanok
\difficulty{4}
Let $s$ be a finite set and let $d$ assign a nonnegative integer depth to each element
of $s$. Then $2^{\max_{i\in s}d(i)}\le 1+\sum_{i\in s}2^{d(i)}$, where the maximum of
the empty set is taken to be $0$.
\end{theorem}

\begin{theorem}[A natural number bounded by a power expression]
\lean{FABL.natCast_le_add_pow_div_pow}
\label{FABL.natCast_le_add_pow_div_pow}
\leanok
\difficulty{5}
For all natural numbers $d$ and $q$, viewing them as real numbers gives $d\le
q+2^d/2^q$.
\end{theorem}

\begin{theorem}[Expected maximum restricted decision tree depth]
\lean{FABL.expect_maxRestrictedDecisionTreeDepth_le_clog_add_two_of_failure_le_quarter}
\label{FABL.expect_maxRestrictedDecisionTreeDepth_le_clog_add_two_of_failure_le_quarter}
\leanok
\uses{FABL.BooleanFunction, FABL.FixedSignCube, FABL.expectRandomRestriction, FABL.expectRandomRestriction_add, FABL.expectRandomRestriction_const, FABL.expectRandomRestriction_const_mul, FABL.expectRandomRestriction_mono, FABL.expect_two_pow_restrictedDecisionTreeDepth_le_two_of_failure_le_quarter, FABL.maxRestrictedDecisionTreeDepth, FABL.natCast_le_add_pow_div_pow, FABL.restrictedDecisionTreeDepth, FABL.sum_expectRandomRestriction, FABL.switchingFailureProbability, FABL.two_pow_sup_le_one_add_sum}
\difficulty{5}
Let $m,n$ be nonnegative integers, let $f_0,\ldots,f_{m-1}$ be Boolean functions on the
$n$-dimensional sign cube, and let $\delta\in[0,1]$. Suppose that for every function in
the family and every nonnegative integer $k$, the probability under a random restriction
with parameter $\delta$ that its restricted decision tree has depth at least $k$ is at
most $(1/4)^k$. Then the expected maximum restricted decision tree depth of the family
is at most $\lceil\log_2(m+1)\rceil+2$.
\end{theorem}

\begin{definition}[Circuit influence logarithmic bound]
\lean{FABL.circuitInfluenceLog}
\label{FABL.circuitInfluenceLog}
\leanok
For a natural number $s$, this quantity is the least natural number $k$ such that
$s+1\le 2^k$, plus $2$.
\end{definition}

\begin{definition}[Circuit influence step]
\lean{FABL.circuitInfluenceStep}
\label{FABL.circuitInfluenceStep}
\leanok
\uses{FABL.circuitInfluenceLog}
For each natural number $s$, define the circuit influence step to be $20$ times the
circuit influence logarithm of $s$.
\end{definition}

\begin{definition}[Quarter-rate switching bound]
\lean{FABL.HasQuarterSwitchingBound}
\statementsource{Analysis-of-Boolean-Functions-by-Ryan-ODonnell}{pdf-ffbe9c7b3306-p100-b005}
\label{FABL.HasQuarterSwitchingBound}
\leanok
\uses{FABL.BooleanFunction, FABL.HasCNFWidthLE, FABL.HasDNFWidthLE, FABL.switchingFailureProbability}
A Boolean function has the quarter-rate switching bound if, for every natural number
$N$, every natural number $w>0$, and every Boolean function $f$ on $N$ coordinates that
has a DNF or a CNF of width at most $w$, the probability that $f$ fails to switch under
restriction parameter $1/(20w)$ after $k$ rounds is at most $(1/4)^k$ for every natural
number $k$.
\end{definition}

\begin{theorem}[Gate width yields a normal-form width bound]
\lean{FABL.CircuitGate.hasDNFWidthLE_or_hasCNFWidthLE}
\proofsource{Analysis-of-Boolean-Functions-by-Ryan-ODonnell}{
  pdf-ffbe9c7b3306-p094-b004,
  pdf-ffbe9c7b3306-p094-b005
}
\label{FABL.CircuitGate.hasDNFWidthLE_or_hasCNFWidthLE}
\leanok
\uses{FABL.BooleanFunction, FABL.CircuitGate, FABL.CircuitGate.HasWidthLE, FABL.HasCNFWidthLE, FABL.HasDNFWidthLE}
\difficulty{4}
Let $n$ be a nonnegative integer, let $f:\{\pm1\}^n\to\{\pm1\}$, let $w$ be a
nonnegative integer, and let $G$ be an AND or OR gate. If $f$ has width at most $w$ in
the normal form associated with $G$, then $f$ has width at most $w$ in DNF or in CNF.
\end{theorem}

\begin{theorem}[Total influence under quarter switching]
\lean{FABL.CircuitTail.totalInfluence_evalFunction_le_of_quarterSwitching}
\label{FABL.CircuitTail.totalInfluence_evalFunction_le_of_quarterSwitching}
\leanok
\uses{FABL.BooleanFunction, FABL.BooleanFunction.toReal, FABL.CircuitGate, FABL.CircuitGate.HasWidthLE, FABL.CircuitGate.dual, FABL.CircuitGate.evalFinsetFunction, FABL.CircuitGate.hasDNFWidthLE_or_hasCNFWidthLE, FABL.CircuitGate.hasWidthLE_evalFinsetFunction, FABL.CircuitGate.hasWidthLE_extendedSignRestriction_max, FABL.CircuitTail, FABL.CircuitTail.evalFunction, FABL.CircuitTail.evalFunction_output, FABL.CircuitTail.extendedSignRestriction_evalFunction, FABL.CircuitTail.internalNodeCount, FABL.CircuitTail.layerCount, FABL.CircuitTail.one_le_layerCount, FABL.CircuitTail.totalInfluence_evalFunction_eq_zero_of_width_zero, FABL.FixedSignCube, FABL.HasQuarterSwitchingBound, FABL.circuitInfluenceLog, FABL.circuitInfluenceStep, FABL.expectRandomRestriction, FABL.expectRandomRestriction_const_mul, FABL.expectRandomRestriction_mono, FABL.expect_maxRestrictedDecisionTreeDepth_le_clog_add_two_of_failure_le_quarter, FABL.expect_totalInfluence_extended_randomRestriction, FABL.extendedSignRestriction, FABL.maxRestrictedDecisionTreeDepth, FABL.switchingFailureProbability, FABL.totalInfluence, FABL.totalInfluence_le_two_mul_of_hasWidthLE}
\difficulty{6}
Assume that every Boolean function with a disjunctive or conjunctive normal form of
width at most a positive integer $w$ has switching failure probability at rate $1/(20w)$
at most $4^{-k}$ for every natural number $k$. Let $T$ be a circuit tail with $m$
inputs, let $g$ be an AND or OR gate, and suppose each input function has a normal form
appropriate to $g$ of width at most $w$. If $N$ is the number of internal gates of $T$
and $N+1\le s$, then the total influence of the real-valued function computed by $T$ is
at most $2w\bigl(20(\lceil\log_2(s+1)\rceil+2)\bigr)^{L-1}$, where $L$ is the number of
layers of $T$.
\end{theorem}

\begin{definition}[Layer-one gate function]
\lean{FABL.DepthCircuit.layer1Function}
\statementsource{Analysis-of-Boolean-Functions-by-Ryan-ODonnell}{pdf-ffbe9c7b3306-p103-b002}
\label{FABL.DepthCircuit.layer1Function}
\leanok
\uses{FABL.BooleanFunction, FABL.DepthCircuit, FABL.DepthCircuit.evalLayer1}
For a depth circuit $C$ on $n$ coordinates and an index $i$ of a layer-one gate, the
associated Boolean function sends each input $x$ to the value produced by gate $i$ in
the first layer of $C$.
\end{definition}

\begin{theorem}[Width bound for a first-layer function]
\lean{FABL.DepthCircuit.layer1Function_hasWidthLE}
\label{FABL.DepthCircuit.layer1Function_hasWidthLE}
\leanok
\uses{FABL.CNFFormula, FABL.CNFFormula.clauseEval, FABL.CNFFormula.eval, FABL.CircuitGate.HasWidthLE, FABL.CircuitGate.dual, FABL.CircuitGate.evalTerm, FABL.DNFFormula, FABL.DNFFormula.eval, FABL.DNFFormula.width_le_of_mem, FABL.DNFTerm.eval, FABL.DNFTerm.width, FABL.DepthCircuit, FABL.DepthCircuit.layer1Function, FABL.DepthCircuit.width}
\difficulty{5}
Let $C$ be a depth circuit on the $n$-dimensional sign cube, and let $i$ index one of
its first-layer gates. If the first-layer gates are AND gates, the function computed by
gate $i$ has a CNF representation of width at most the width of $C$; if they are OR
gates, it has a DNF representation of width at most the width of $C$.
\end{theorem}

\begin{theorem}[First-layer evaluation gives circuit function]
\lean{FABL.DepthCircuit.evalFunction_layer1}
\label{FABL.DepthCircuit.evalFunction_layer1}
\leanok
\uses{FABL.CircuitGate.dual, FABL.CircuitTail.evalFunction, FABL.DepthCircuit, FABL.DepthCircuit.layer1Function, FABL.DepthCircuit.toBooleanFunction}
\difficulty{1}
For every natural number $n$ and every circuit on $n$ variables, evaluating its later
layers on the functions computed by its first-layer gates gives exactly the function
computed by the circuit.
\end{theorem}

\begin{definition}[Circuit with depth, width, and tail size bounds]
\lean{FABL.DepthCircuit.HasDepthWidthTailSizeCircuit}
\statementsource{Analysis-of-Boolean-Functions-by-Ryan-ODonnell}{
  pdf-ffbe9c7b3306-p103-b002,
  pdf-ffbe9c7b3306-p104-b001
}
\label{FABL.DepthCircuit.HasDepthWidthTailSizeCircuit}
\leanok
\uses{FABL.BooleanFunction, FABL.CircuitTail.internalNodeCount, FABL.DepthCircuit, FABL.DepthCircuit.depth, FABL.DepthCircuit.toBooleanFunction, FABL.DepthCircuit.width}
A Boolean function $f$ on $n$ coordinates has a depth-$d$ circuit of width at most $w$
and tail size at most $s'$ if there is a depth circuit $C$ computing $f$ such that its
depth is exactly $d$, its width is at most $w$, and the number of internal nodes in its
tail, plus one, is at most $s'$.
\end{definition}

\begin{theorem}[Total influence under a switching bound]
\lean{FABL.DepthCircuit.exercise4_20_of_quarterSwitching}
\label{FABL.DepthCircuit.exercise4_20_of_quarterSwitching}
\leanok
\uses{FABL.BooleanFunction, FABL.BooleanFunction.toReal, FABL.CircuitGate.HasWidthLE, FABL.CircuitGate.dual, FABL.CircuitGate.hasWidthLE_mono, FABL.CircuitTail.layerCount, FABL.CircuitTail.one_le_layerCount, FABL.CircuitTail.totalInfluence_evalFunction_le_of_quarterSwitching, FABL.DepthCircuit.HasDepthWidthTailSizeCircuit, FABL.DepthCircuit.depth, FABL.DepthCircuit.evalFunction_layer1, FABL.DepthCircuit.layer1Function, FABL.DepthCircuit.layer1Function_hasWidthLE, FABL.HasQuarterSwitchingBound, FABL.circuitInfluenceStep, FABL.totalInfluence}
\difficulty{4}
Assume that for every dimension $N$, every positive width $w$, every Boolean function
$f:\{−1,1\}^N\to\{−1,1\}$ representable by a DNF or CNF of width at most $w$, and every
natural number $k$, the probability that a random restriction with each coordinate fixed
with probability $1/(20w)$ leaves $f$ with decision-tree depth at least $k$ is at most
$(1/4)^k$. If $f:\{−1,1\}^n\to\{−1,1\}$ is computed by a layered circuit of depth $d$,
first-layer width at most $w$, and tail size at most $s'$, where tail size is the number
of non-output gates in layers after the first, plus one, then $\operatorname{Inf}(f)\le
2w\bigl(20(\lceil\log_2(s'+1)\rceil+2)\bigr)^{d-2}$.
\end{theorem}

\begin{theorem}[Hastad quarter switching bound]
\lean{FABL.hastadQuarterSwitchingBound}
\proofsource{lmn-constant-depth-fourier}{pdf-63759bd909d5-p005-b004}
\statementsource{Analysis-of-Boolean-Functions-by-Ryan-ODonnell}{pdf-ffbe9c7b3306-p100-b005}
\label{FABL.hastadQuarterSwitchingBound}
\leanok
\uses{FABL.HasQuarterSwitchingBound, FABL.hastadSwitchingLemma, FABL.switchingFailureProbability}
\difficulty{4}
% dataset-pipeline: stub -- prose pending build_informal_statements.py
Statement pending regeneration.
\end{theorem}

\begin{definition}[Completing a restricted sign vector]
\lean{FABL.CoordRestriction.completeSign}
\label{FABL.CoordRestriction.completeSign}
\leanok
\uses{FABL.CoordRestriction}
Given a coordinate restriction on $n$ coordinates and a sign vector $x\in\{-1,1\}^n$,
define its completion by retaining $x_i$ at each free coordinate, setting coordinates
fixed to $1$ equal to $1$, and setting coordinates fixed to $-1$ equal to $-1$.
\end{definition}

\begin{definition}[Coordinate restriction of a sign function]
\lean{FABL.CoordRestriction.restrictSign}
\statementsource{Analysis-of-Boolean-Functions-by-Ryan-ODonnell}{
  pdf-ffbe9c7b3306-p098-b003,
  pdf-ffbe9c7b3306-p098-b004
}
\label{FABL.CoordRestriction.restrictSign}
\leanok
\uses{FABL.BooleanFunction, FABL.CoordRestriction, FABL.CoordRestriction.completeSign}
Given a Boolean function $f$ on $n$ coordinates and a coordinate restriction $
ho$ that
either leaves each coordinate free or fixes it to $1$ or $-1$, the restricted function
sends an input $x$ to the value of $f$ at the completion of $x$ specified by $
ho$.
\end{definition}

\begin{theorem}[Completion from a fixed sign assignment]
\lean{FABL.CoordRestriction.completeSign_coordRestrictionOf}
\label{FABL.CoordRestriction.completeSign_coordRestrictionOf}
\leanok
\uses{FABL.CoordRestriction.completeSign, FABL.FixedSignCube, FABL.combineSignCube, FABL.combineSignCube_apply_fixed, FABL.combineSignCube_apply_free, FABL.coordRestrictionOf}
\difficulty{4}
Let $n$ be a natural number, let $J$ be a set of coordinates among the $n$ coordinates,
let $z$ assign a sign to every coordinate outside $J$, and let $x$ be a sign vector of
length $n$. The sign vector obtained by completing the restriction that leaves $J$ free
and fixes every coordinate outside $J$ according to $z$, using $x$ on the free
coordinates, is the vector formed from the restriction of $x$ to $J$ and the assignment
$z$ outside $J$.
\end{theorem}

\begin{theorem}[Restriction from fixed coordinates]
\lean{FABL.CoordRestriction.restrictSign_coordRestrictionOf}
\label{FABL.CoordRestriction.restrictSign_coordRestrictionOf}
\leanok
\uses{FABL.BooleanFunction, FABL.CoordRestriction.completeSign_coordRestrictionOf, FABL.CoordRestriction.restrictSign, FABL.FixedSignCube, FABL.coordRestrictionOf, FABL.extendedSignRestriction}
\difficulty{3}
Let $n$ be a nonnegative integer, let $f$ be a function from the $n$-dimensional sign
cube to $\{\pm1\}$, let $J$ be a set of coordinates, and let $z$ assign a sign to each
coordinate outside $J$. Restricting $f$ according to the coordinates fixed by $z$ and
leaving the coordinates in $J$ free is the same function as evaluating $f$ after
combining an arbitrary assignment on $J$ with $z$ on the other coordinates.
\end{theorem}

\begin{theorem}[Completion commutes with sign negation]
\lean{FABL.CoordRestriction.completeSign_negateAssignment}
\label{FABL.CoordRestriction.completeSign_negateAssignment}
\leanok
\uses{FABL.CoordRestriction, FABL.CoordRestriction.completeSign, FABL.CoordRestriction.negateAssignment, FABL.CoordRestriction.negateState}
\difficulty{4}
Let $n$ be a natural number, let $\rho$ be a coordinate restriction on $n$ coordinates,
and let $x$ be a sign vector of length $n$. Completing the restriction obtained by
reversing every fixed sign in $\rho$, using the coordinatewise negation of $x$, gives
the coordinatewise negation of the completion of $\rho$ using $x$.
\end{theorem}

\begin{theorem}[Duality and sign restriction]
\lean{FABL.CoordRestriction.restrictSign_booleanDual}
\label{FABL.CoordRestriction.restrictSign_booleanDual}
\leanok
\uses{FABL.BooleanFunction, FABL.CNFFormula.booleanDual, FABL.CoordRestriction, FABL.CoordRestriction.completeSign_negateAssignment, FABL.CoordRestriction.negateAssignment, FABL.CoordRestriction.restrictSign}
\difficulty{3}
Let $n$ be a nonnegative integer, let $f$ be a function from the $n$-dimensional sign
cube to $\\{\pm1\}$, and let $\rho$ assign to each coordinate one of three states: free,
fixed to $1$, or fixed to $-1$. Restricting the Boolean dual of $f$ using $\rho$ gives
the Boolean dual of $f$ restricted using the coordinatewise negation of $\rho$.
\end{theorem}

\begin{definition}[Free literals of a restricted term]
\lean{FABL.DNFTerm.freeRestriction}
\statementsource{Analysis-of-Boolean-Functions-by-Ryan-ODonnell}{
  pdf-ffbe9c7b3306-p098-b003,
  pdf-ffbe9c7b3306-p098-b004
}
\label{FABL.DNFTerm.freeRestriction}
\leanok
\uses{FABL.CoordRestriction, FABL.DNFTerm}
Given a DNF term $T$ on $n$ coordinates and a coordinate restriction $
ho$, the free
restriction of $T$ is the term consisting of precisely those literals of $T$ whose
coordinates are left free by $
ho$.
\end{definition}

\begin{definition}[Formula induced by a restricted term]
\lean{FABL.DNFTerm.restrictionFormula}
\statementsource{Analysis-of-Boolean-Functions-by-Ryan-ODonnell}{
  pdf-ffbe9c7b3306-p098-b003,
  pdf-ffbe9c7b3306-p098-b004
}
\label{FABL.DNFTerm.restrictionFormula}
\leanok
\uses{FABL.CoordRestriction, FABL.DNFFormula, FABL.DNFFormula.empty, FABL.DNFTerm, FABL.DNFTerm.freeRestriction, FABL.DNFTerm.notFalsified}
Given a DNF term $T$ on $n$ coordinates and a coordinate restriction $
ho$, its
restriction formula is the one-term DNF consisting of the free restriction of $T$ if the
restriction does not falsify $T$, and is the empty DNF otherwise.
\end{definition}

\begin{theorem}[Width after restricting a term]
\lean{FABL.DNFTerm.width_restrictionFormula}
\label{FABL.DNFTerm.width_restrictionFormula}
\leanok
\uses{FABL.CoordRestriction, FABL.DNFFormula.empty, FABL.DNFFormula.width, FABL.DNFTerm, FABL.DNFTerm.freeRestriction, FABL.DNFTerm.notFalsified, FABL.DNFTerm.restrictedWidthOf, FABL.DNFTerm.restrictionFormula, FABL.DNFTerm.width, FABL.Literal.isFalsified}
\difficulty{4}
Let $T$ be a DNF term on $n$ coordinates, and assign each coordinate either the
restriction free, fixed to $1$, or fixed to $-1$. The width of the formula obtained by
restricting $T$ equals $0$ if any literal of $T$ is falsified by a fixed coordinate;
otherwise it equals the number of literals of $T$ whose coordinates are free.
\end{theorem}

\begin{theorem}[Signs determined by the negative-sign test]
\lean{FABL.DNFTerm.sign_eq_of_neg_one_iff}
\label{FABL.DNFTerm.sign_eq_of_neg_one_iff}
\leanok
\uses{FABL.Sign}
\difficulty{4}
For signs $a,b\in\{ -1,1\}$, if $a=-1$ if and only if $b=-1$, then $a=b$.
\end{theorem}

\begin{theorem}[Restriction formula represents the restricted term]
\lean{FABL.DNFTerm.restrictionFormula_toBooleanFunction}
\label{FABL.DNFTerm.restrictionFormula_toBooleanFunction}
\leanok
\uses{FABL.CoordRestriction, FABL.CoordRestriction.completeSign, FABL.CoordRestriction.restrictSign, FABL.DNFFormula.eval_eq_neg_one_iff, FABL.DNFFormula.toBooleanFunction, FABL.DNFTerm, FABL.DNFTerm.eval, FABL.DNFTerm.eval_eq_neg_one_iff, FABL.DNFTerm.freeRestriction, FABL.DNFTerm.notFalsified, FABL.DNFTerm.notFalsified_iff, FABL.DNFTerm.restrictionFormula, FABL.DNFTerm.sign_eq_of_neg_one_iff, FABL.Literal.isFalsified, FABL.Sign}
\difficulty{5}
Let $T$ be a DNF term on $n$ coordinates, and let each coordinate be free, fixed to $1$,
or fixed to $-1$. The Boolean function represented by the restriction formula of $T$
agrees on every sign vector with the function obtained by evaluating $T$ after
completing the restriction using that vector on the free coordinates.
\end{theorem}

\begin{theorem}[DNF width after restriction]
\lean{FABL.DNFTerm.hasDNFWidthLE_restrictSign}
\label{FABL.DNFTerm.hasDNFWidthLE_restrictSign}
\leanok
\uses{FABL.CoordRestriction, FABL.CoordRestriction.restrictSign, FABL.DNFTerm, FABL.DNFTerm.eval, FABL.DNFTerm.restrictedWidthOf, FABL.DNFTerm.restrictionFormula, FABL.DNFTerm.restrictionFormula_toBooleanFunction, FABL.DNFTerm.width_restrictionFormula, FABL.HasDNFWidthLE}
\difficulty{1}
Let $T$ be a conjunction of literals on $n$ coordinates, and let $\rho$ be a coordinate
restriction. The function obtained by restricting the function represented by $T$ has a
disjunctive normal form representation of width at most the restricted width of $T$
under $\rho$.
\end{theorem}

\begin{theorem}[Clause evaluation and Boolean duality]
\lean{FABL.CNFFormula.clauseEval_eq_booleanDual_eval}
\proofsource{Analysis-of-Boolean-Functions-by-Ryan-ODonnell}{
  pdf-ffbe9c7b3306-p094-b004,
  pdf-ffbe9c7b3306-p094-b005
}
\label{FABL.CNFFormula.clauseEval_eq_booleanDual_eval}
\leanok
\uses{FABL.CNFFormula.booleanDual, FABL.CNFFormula.clauseEval, FABL.CNFFormula.clauseEval_neg_iff_termEval, FABL.DNFTerm, FABL.DNFTerm.eval}
\difficulty{4}
Let $T$ be a term on the sign cube of dimension $n$. As functions on that cube, the
evaluation of the clause formed from $T$ equals the Boolean dual of the evaluation of
$T$; equivalently, at every assignment $x$, the clause evaluates to the negative of the
term's evaluation at the coordinatewise negation of $x$.
\end{theorem}

\begin{theorem}[Boolean dual]
\lean{FABL.HasDNFWidthLE.booleanDual}
\statementsource{Analysis-of-Boolean-Functions-by-Ryan-ODonnell}{
  pdf-ffbe9c7b3306-p094-b004,
  pdf-ffbe9c7b3306-p094-b005
}
\label{FABL.HasDNFWidthLE.booleanDual}
\leanok
\uses{FABL.BooleanFunction, FABL.CNFFormula, FABL.CNFFormula.booleanDual, FABL.CNFFormula.switchAndOr, FABL.CNFFormula.switchAndOr_toBooleanFunction, FABL.CNFFormula.toBooleanFunction, FABL.CNFFormula.width, FABL.DNFFormula.toBooleanFunction, FABL.HasCNFWidthLE, FABL.HasDNFWidthLE, FABL.booleanDual_injective}
\difficulty{4}
% dataset-pipeline: stub -- prose pending build_informal_statements.py
Statement pending regeneration.
\end{theorem}

\begin{theorem}[CNF width after restricting a clause]
\lean{FABL.DNFTerm.hasCNFWidthLE_restrictSign_clause}
\label{FABL.DNFTerm.hasCNFWidthLE_restrictSign_clause}
\leanok
\uses{FABL.CNFFormula.clauseEval, FABL.CNFFormula.clauseEval_eq_booleanDual_eval, FABL.CoordRestriction, FABL.CoordRestriction.negateAssignment, FABL.CoordRestriction.restrictSign, FABL.CoordRestriction.restrictSign_booleanDual, FABL.DNFTerm, FABL.DNFTerm.hasDNFWidthLE_restrictSign, FABL.DNFTerm.restrictedWidthOf, FABL.HasCNFWidthLE, FABL.HasDNFWidthLE.booleanDual}
\difficulty{3}
Let $T$ be a conjunction of literals on $n$ coordinates, and let $\rho$ be a coordinate
restriction. The clause obtained by restricting the clause represented by $T$ has a
conjunctive normal form representation of width at most the restricted width of $T$,
computed using the restriction with its fixed signs negated.
\end{theorem}

\begin{definition}[Restricted gate width]
\lean{FABL.DNFTerm.restrictedGateWidth}
\label{FABL.DNFTerm.restrictedGateWidth}
\leanok
\uses{FABL.CircuitGate, FABL.CoordRestriction, FABL.CoordRestriction.negateAssignment, FABL.DNFTerm, FABL.DNFTerm.restrictedWidthOf}
For a DNF term $T$, a coordinate restriction $\rho$, and a gate type, define the
restricted gate width to be the restricted width of $T$ under $\rho$ when the gate is
AND, and the restricted width under the negated assignment restriction when the gate is
OR.
\end{definition}

\begin{theorem}[Width of a restricted gate term]
\lean{FABL.CircuitGate.hasWidthLE_restrictSign_evalTerm}
\label{FABL.CircuitGate.hasWidthLE_restrictSign_evalTerm}
\leanok
\uses{FABL.CNFFormula.clauseEval, FABL.CircuitGate, FABL.CircuitGate.HasWidthLE, FABL.CircuitGate.dual, FABL.CircuitGate.evalTerm, FABL.CoordRestriction, FABL.CoordRestriction.negateAssignment, FABL.CoordRestriction.restrictSign, FABL.DNFTerm, FABL.DNFTerm.eval, FABL.DNFTerm.hasCNFWidthLE_restrictSign_clause, FABL.DNFTerm.hasDNFWidthLE_restrictSign, FABL.DNFTerm.restrictedGateWidth, FABL.DNFTerm.restrictedWidthOf, FABL.HasCNFWidthLE, FABL.HasDNFWidthLE}
\difficulty{4}
Let $n$ be a natural number, choose an AND or OR gate, a term of literals on $n$
coordinates, and a coordinate restriction assigning each coordinate to be free, fixed at
$1$, or fixed at $-1$. Apply the restriction to the function computed by the chosen gate
on the term. The resulting function has a representation of width at most the term's
restricted gate width, using the dual gate (OR for AND, and AND for OR).
\end{theorem}

\begin{definition}[Maximum first-layer restricted width]
\lean{FABL.DepthCircuit.maxLayer1RestrictedWidth}
\label{FABL.DepthCircuit.maxLayer1RestrictedWidth}
\leanok
\uses{FABL.CoordRestriction, FABL.DNFTerm.restrictedGateWidth, FABL.DepthCircuit}
For a depth circuit $C$ on $n$ coordinates and a coordinate restriction $\rho$, define
its maximum first-layer restricted width as the supremum, over all first-layer gates, of
each gate's restricted gate width under $\rho$.
\end{definition}

\begin{theorem}[Restricted width bounded by the maximum]
\lean{FABL.DepthCircuit.layer1RestrictedWidth_le_max}
\label{FABL.DepthCircuit.layer1RestrictedWidth_le_max}
\leanok
\uses{FABL.CoordRestriction, FABL.DNFTerm.restrictedGateWidth, FABL.DepthCircuit, FABL.DepthCircuit.maxLayer1RestrictedWidth}
\difficulty{1}
Let $C$ be a depth circuit, let $\rho$ assign each input coordinate one of the states
free, fixed to $1$, or fixed to $-1$, and let $i$ index a first-layer gate. The
restricted gate width of gate $i$ under $\rho$ is at most the maximum restricted width
of the first-layer gates under $\rho$.
\end{theorem}

\begin{theorem}[Width bound after a sign restriction]
\lean{FABL.DepthCircuit.layer1Function_hasWidthLE_extendedSignRestriction_max}
\label{FABL.DepthCircuit.layer1Function_hasWidthLE_extendedSignRestriction_max}
\leanok
\uses{FABL.CircuitGate.HasWidthLE, FABL.CircuitGate.dual, FABL.CircuitGate.evalTerm, FABL.CircuitGate.hasWidthLE_mono, FABL.CircuitGate.hasWidthLE_restrictSign_evalTerm, FABL.CoordRestriction.restrictSign_coordRestrictionOf, FABL.DepthCircuit, FABL.DepthCircuit.layer1Function, FABL.DepthCircuit.layer1RestrictedWidth_le_max, FABL.DepthCircuit.maxLayer1RestrictedWidth, FABL.FixedSignCube, FABL.coordRestrictionOf, FABL.extendedSignRestriction}
\difficulty{3}
Let $C$ be a depth circuit on the $n$-dimensional sign cube, let $J$ be a set of
coordinates, let $z$ assign a sign to every coordinate outside $J$, and let $i$ index a
first-layer gate. Fix the coordinates outside $J$ according to $z$ and leave the
coordinates in $J$ free. The resulting restricted function of gate $i$ has, if the
first-layer gates are AND gates, a CNF representation, and if they are OR gates, a DNF
representation, whose width is at most the maximum restricted width among the
first-layer gates under this coordinate restriction.
\end{theorem}

\begin{theorem}[Restriction assignment weight at half]
\lean{FABL.restrictionAssignmentWeightAt_half}
\label{FABL.restrictionAssignmentWeightAt_half}
\leanok
\uses{FABL.CoordRestriction, FABL.coordRestrictionWeight, FABL.coordRestrictionWeightAt, FABL.restrictionAssignmentWeight, FABL.restrictionAssignmentWeightAt}
\difficulty{4}
For any natural number $n$ and any assignment of one of the three states “free,” “fixed
to $1$,” or “fixed to $-1$” to each of the $n$ coordinates, its product weight at
parameter $rac12$ equals its ordinary product weight, where the ordinary coordinate
weights are $rac12$ for a free coordinate and $rac14$ for either fixed coordinate.
\end{theorem}

\begin{theorem}[Negation invariance of assignment weight]
\lean{FABL.restrictionAssignmentWeight_negateAssignment}
\label{FABL.restrictionAssignmentWeight_negateAssignment}
\leanok
\uses{FABL.CoordRestriction, FABL.CoordRestriction.negateAssignment, FABL.CoordRestriction.negateState, FABL.coordRestrictionWeight, FABL.restrictionAssignmentWeight}
\difficulty{4}
Let $n$ be a natural number and assign to each of the $n$ coordinates one of the states
“free,” “fixed to $1$,” or “fixed to $-1$.” Exchanging the two fixed states in every
coordinate leaves the product weight unchanged, where the coordinate weights are
$rac12$ for a free state and $rac14$ for either fixed state.
\end{theorem}

\begin{definition}[Blocked restricted gate width]
\lean{FABL.DNFTerm.blockedRestrictedGateWidth}
\label{FABL.DNFTerm.blockedRestrictedGateWidth}
\leanok
\uses{FABL.CircuitGate, FABL.CoordRestriction, FABL.DNFTerm, FABL.DNFTerm.restrictedGateWidth}
For a DNF term $T$, a gate type, and a coordinate restriction $\rho$, define the blocked
restricted gate width to be the restricted gate width divided by $5$, with division in
the natural numbers.
\end{definition}

\begin{theorem}[Exponential power comparison]
\lean{FABL.three_quarters_pow_five_mul_le_quarter_pow}
\label{FABL.three_quarters_pow_five_mul_le_quarter_pow}
\leanok
\difficulty{3}
For every natural number $k$, $(3/4)^{5k}\le (1/4)^k$.
\end{theorem}

\begin{theorem}[Tail bound for blocked restricted gate width]
\lean{FABL.DNFTerm.expect_blockedRestrictedGateWidth_failure_le_quarter}
\proofsource{Analysis-of-Boolean-Functions-by-Ryan-ODonnell}{
  pdf-ffbe9c7b3306-p099-b005,
  pdf-ffbe9c7b3306-p099-b006
}
\label{FABL.DNFTerm.expect_blockedRestrictedGateWidth_failure_le_quarter}
\leanok
\uses{FABL.CircuitGate, FABL.CoordRestriction, FABL.CoordRestriction.negateAssignment, FABL.DNFTerm, FABL.DNFTerm.blockedRestrictedGateWidth, FABL.DNFTerm.restrictedGateWidth, FABL.DNFTerm.restrictedWidthOf, FABL.coordRestrictionNegationEquiv, FABL.coordRestrictionOf, FABL.coordRestrictionOf_restrictionAtomOfAssignment, FABL.expectRandomRestriction, FABL.expectRandomRestriction_eq_sum_assignment, FABL.restrictedWidth_ge_probability_le, FABL.restrictionAssignmentWeight, FABL.restrictionAssignmentWeightAt_half, FABL.restrictionAssignmentWeight_nonneg, FABL.three_quarters_pow_five_mul_le_quarter_pow}
\difficulty{5}
Let $T$ be a conjunction of literals on $n$ coordinates, let $g$ be an AND or OR gate,
and let $k$ be a natural number. Choose a subset $J\subseteq\{1,\ldots,n\}$ by including
each coordinate independently with probability $1/2$, then choose uniformly a sign
assignment $z$ on the coordinates outside $J$. Form the restriction that leaves $J$ free
and fixes the other coordinates according to $z$. The probability that the restricted
gate width of $T$, divided by $5$ and rounded down, is at least $k$ is at most
$(1/4)^k$.
\end{theorem}

\begin{theorem}[Exponential depth bound from tail estimates]
\lean{FABL.expect_two_pow_le_two_of_tail_le_quarter}
\label{FABL.expect_two_pow_le_two_of_tail_le_quarter}
\leanok
\uses{FABL.FixedSignCube, FABL.expectRandomRestriction, FABL.expectRandomRestriction_add, FABL.expectRandomRestriction_const, FABL.expectRandomRestriction_const_mul, FABL.expectRandomRestriction_sum, FABL.two_pow_eq_one_add_sum_tailIndicators, FABL.two_pow_mul_quarter_pow_succ}
\difficulty{4}
Let $n$ be a nonnegative integer and let $\delta$ be any real number. For each subset
$J\subseteq\{1,\ldots,n\}$ and each assignment $z$ of signs outside $J$, let $d(J,z)$ be
a nonnegative integer, and let $N$ be a nonnegative integer such that $d(J,z)\le N$ for
every $J,z$. Use the random-restriction average that weights $J$ by
$\delta^{|J|}(1-\delta)^{n-|J|}$ and averages uniformly over $z$ outside $J$. If, for
every nonnegative integer $k$, the average of the indicator of $k\le d(J,z)$ is at most
$(1/4)^k$, then the average of $2^{d(J,z)}$ is at most $2$.
\end{theorem}

\begin{theorem}[Expect two pow blocked restricted gate width le two]
\lean{FABL.DNFTerm.expect_two_pow_blockedRestrictedGateWidth_le_two}
\statementsource{Analysis-of-Boolean-Functions-by-Ryan-ODonnell}{
  pdf-ffbe9c7b3306-p099-b005,
  pdf-ffbe9c7b3306-p099-b006
}
\label{FABL.DNFTerm.expect_two_pow_blockedRestrictedGateWidth_le_two}
\leanok
\uses{FABL.CircuitGate, FABL.DNFTerm, FABL.DNFTerm.blockedRestrictedGateWidth, FABL.DNFTerm.expect_blockedRestrictedGateWidth_failure_le_quarter, FABL.DNFTerm.restrictedWidthOf_le_width, FABL.DNFTerm.width_le_dimension, FABL.coordRestrictionOf, FABL.expectRandomRestriction, FABL.expect_two_pow_le_two_of_tail_le_quarter}
\difficulty{4}
% dataset-pipeline: stub -- prose pending build_informal_statements.py
Statement pending regeneration.
\end{theorem}

\begin{theorem}[Expected maximum depth bound]
\lean{FABL.expect_sup_le_clog_add_two_of_twoPow_le_two}
\label{FABL.expect_sup_le_clog_add_two_of_twoPow_le_two}
\leanok
\uses{FABL.FixedSignCube, FABL.expectRandomRestriction, FABL.expectRandomRestriction_add, FABL.expectRandomRestriction_const, FABL.expectRandomRestriction_const_mul, FABL.expectRandomRestriction_mono, FABL.natCast_le_add_pow_div_pow, FABL.sum_expectRandomRestriction, FABL.two_pow_sup_le_one_add_sum}
\difficulty{5}
Let $m,n$ be natural numbers. For each $i\in\{1,\ldots,m\}$, let the depth be a
natural-number-valued function of a subset $J\subseteq\{1,\ldots,n\}$ and a sign
assignment on the coordinates outside $J$. Suppose $0\leq\delta\leq1$, and suppose that
for every $i$, the expected value of $2^{\mathrm{depth}_i}$ under the random restriction
with parameter $\delta$ is at most $2$. Then the expected maximum of these depths under
the same random restriction is at most $\lceil\log_2(m+1)\rceil+2$.
\end{theorem}

\begin{definition}[Maximum first-layer blocked width]
\lean{FABL.DepthCircuit.maxLayer1BlockedWidth}
\label{FABL.DepthCircuit.maxLayer1BlockedWidth}
\leanok
\uses{FABL.CoordRestriction, FABL.DNFTerm.blockedRestrictedGateWidth, FABL.DepthCircuit}
For a depth circuit $C$ on $n$ coordinates and a coordinate restriction $\rho$, define
its maximum first-layer blocked width as the supremum, over all first-layer gates, of
each gate's blocked restricted width under $\rho$.
\end{definition}

\begin{theorem}[Restricted width and blocked width]
\lean{FABL.DepthCircuit.maxLayer1RestrictedWidth_le_blocked}
\label{FABL.DepthCircuit.maxLayer1RestrictedWidth_le_blocked}
\leanok
\uses{FABL.CoordRestriction, FABL.DNFTerm.blockedRestrictedGateWidth, FABL.DepthCircuit, FABL.DepthCircuit.maxLayer1BlockedWidth, FABL.DepthCircuit.maxLayer1RestrictedWidth}
\difficulty{3}
Let $C$ be a depth circuit and assign to each input coordinate one of three states:
free, fixed to $1$, or fixed to $-1$. The maximum restricted width among the first-layer
gates is at most five times the maximum blocked restricted width, plus four.
\end{theorem}

\begin{definition}[Initial restriction width]
\lean{FABL.DepthCircuit.initialRestrictionWidth}
\label{FABL.DepthCircuit.initialRestrictionWidth}
\leanok
\uses{FABL.circuitInfluenceLog}
For each natural number $s$, define the initial restriction width to be five times the
circuit influence logarithm of $s$, plus $4$.
\end{definition}

\begin{theorem}[Expect max layer1 restricted width le]
\lean{FABL.DepthCircuit.expect_maxLayer1RestrictedWidth_le}
\statementsource{Analysis-of-Boolean-Functions-by-Ryan-ODonnell}{
  pdf-ffbe9c7b3306-p099-b005,
  pdf-ffbe9c7b3306-p099-b006,
  pdf-ffbe9c7b3306-p100-b001
}
\label{FABL.DepthCircuit.expect_maxLayer1RestrictedWidth_le}
\leanok
\uses{FABL.DNFTerm.blockedRestrictedGateWidth, FABL.DNFTerm.expect_two_pow_blockedRestrictedGateWidth_le_two, FABL.DepthCircuit, FABL.DepthCircuit.initialRestrictionWidth, FABL.DepthCircuit.maxLayer1BlockedWidth, FABL.DepthCircuit.maxLayer1RestrictedWidth, FABL.DepthCircuit.maxLayer1RestrictedWidth_le_blocked, FABL.circuitInfluenceLog, FABL.coordRestrictionOf, FABL.expectRandomRestriction, FABL.expectRandomRestriction_add, FABL.expectRandomRestriction_const, FABL.expectRandomRestriction_const_mul, FABL.expectRandomRestriction_mono, FABL.expect_sup_le_clog_add_two_of_twoPow_le_two}
\difficulty{4}
% dataset-pipeline: stub -- prose pending build_informal_statements.py
Statement pending regeneration.
\end{theorem}

\begin{theorem}[Monotonicity of initial restriction width]
\lean{FABL.DepthCircuit.initialRestrictionWidth_mono}
\label{FABL.DepthCircuit.initialRestrictionWidth_mono}
\leanok
\uses{FABL.DepthCircuit.initialRestrictionWidth, FABL.circuitInfluenceLog}
\difficulty{3}
For natural numbers $a$ and $b$ with $a \le b$, the initial restriction width at $a$ is
at most the initial restriction width at $b$.
\end{theorem}

\begin{theorem}[Restriction-width bound]
\lean{FABL.DepthCircuit.initialRestrictionWidth_le_step_succ}
\label{FABL.DepthCircuit.initialRestrictionWidth_le_step_succ}
\leanok
\uses{FABL.DepthCircuit.initialRestrictionWidth, FABL.circuitInfluenceLog, FABL.circuitInfluenceStep}
\difficulty{4}
For every nonnegative integer $s$, let $L(s)=\lceil\log_2(s+1)\rceil$. Then
$5L(s)+14\leq 20\bigl(L(s+1)+2\bigr)$.
\end{theorem}

\begin{definition}[Circuit with bounded depth and size]
\lean{FABL.DepthCircuit.HasDepthSizeCircuit}
\statementsource{Analysis-of-Boolean-Functions-by-Ryan-ODonnell}{
  pdf-ffbe9c7b3306-p103-b002,
  pdf-ffbe9c7b3306-p104-b001
}
\label{FABL.DepthCircuit.HasDepthSizeCircuit}
\leanok
\uses{FABL.BooleanFunction, FABL.DepthCircuit.HasDepthCircuit}
A Boolean function $f$ on $n$ coordinates has a circuit of depth $d$ and size at most
$s$ if there is a width $w$ for which $f$ has a depth-$d$ circuit of size at most $s$
and width $w$.
\end{definition}

\begin{theorem}[Exercise4 20b of quarter switching]
\lean{FABL.DepthCircuit.exercise4_20b_of_quarterSwitching}
\statementsource{Analysis-of-Boolean-Functions-by-Ryan-ODonnell}{
  pdf-ffbe9c7b3306-p106-b007,
  pdf-ffbe9c7b3306-p100-b005
}
\label{FABL.DepthCircuit.exercise4_20b_of_quarterSwitching}
\leanok
\uses{FABL.BooleanFunction, FABL.BooleanFunction.toReal, FABL.CircuitGate.dual, FABL.CircuitTail.extendedSignRestriction_evalFunction, FABL.CircuitTail.internalNodeCount, FABL.CircuitTail.layerCount, FABL.CircuitTail.one_le_layerCount, FABL.CircuitTail.totalInfluence_evalFunction_le_of_quarterSwitching, FABL.DepthCircuit.HasDepthSizeCircuit, FABL.DepthCircuit.depth, FABL.DepthCircuit.evalFunction_layer1, FABL.DepthCircuit.expect_maxLayer1RestrictedWidth_le, FABL.DepthCircuit.initialRestrictionWidth, FABL.DepthCircuit.initialRestrictionWidth_mono, FABL.DepthCircuit.layer1Function, FABL.DepthCircuit.layer1Function_hasWidthLE_extendedSignRestriction_max, FABL.DepthCircuit.maxLayer1RestrictedWidth, FABL.DepthCircuit.size, FABL.DepthCircuit.toBooleanFunction, FABL.FixedSignCube, FABL.HasQuarterSwitchingBound, FABL.circuitInfluenceLog, FABL.circuitInfluenceStep, FABL.coordRestrictionOf, FABL.expectRandomRestriction, FABL.expectRandomRestriction_const_mul, FABL.expectRandomRestriction_mono, FABL.expect_totalInfluence_extended_randomRestriction, FABL.extendedSignRestriction, FABL.totalInfluence}
\difficulty{5}
% dataset-pipeline: stub -- prose pending build_informal_statements.py
Statement pending regeneration.
\end{theorem}

\begin{theorem}[Influence bound from quarter switching]
\lean{FABL.DepthCircuit.theorem4_30_of_quarterSwitching}
\proofsource{Analysis-of-Boolean-Functions-by-Ryan-ODonnell}{
  pdf-ffbe9c7b3306-p106-b007,
  pdf-ffbe9c7b3306-p110-b004
}
\label{FABL.DepthCircuit.theorem4_30_of_quarterSwitching}
\leanok
\uses{FABL.BooleanFunction, FABL.BooleanFunction.toReal, FABL.DepthCircuit.HasDepthSizeCircuit, FABL.DepthCircuit.depth_ge_two, FABL.DepthCircuit.exercise4_20b_of_quarterSwitching, FABL.DepthCircuit.initialRestrictionWidth, FABL.DepthCircuit.initialRestrictionWidth_le_step_succ, FABL.HasQuarterSwitchingBound, FABL.circuitInfluenceLog, FABL.circuitInfluenceStep, FABL.totalInfluence}
\difficulty{4}
Assume every Boolean function on a finite-dimensional sign cube that has a DNF or CNF
representation of positive width at most $w$ has probability at most $(1/4)^k$ of having
restricted decision-tree depth at least $k$, under a random restriction that leaves each
coordinate free with probability $1/(20w)$. If a Boolean function $f$ on the
$n$-dimensional sign cube is computed by a depth-$d$ circuit of size at most $s$ for
some width, then its total influence is at most
$4\bigl(20(\lceil\log_2(s+1)\rceil+2)\bigr)^{d-1}$.
\end{theorem}

\begin{theorem}[Total influence of a bounded circuit]
\lean{FABL.DepthCircuit.exercise4_20}
\label{FABL.DepthCircuit.exercise4_20}
\leanok
\uses{FABL.BooleanFunction, FABL.BooleanFunction.toReal, FABL.DepthCircuit.HasDepthWidthTailSizeCircuit, FABL.DepthCircuit.exercise4_20_of_quarterSwitching, FABL.circuitInfluenceStep, FABL.hastadQuarterSwitchingBound, FABL.totalInfluence}
\difficulty{1}
Let $f:\{−1,1\}^n\to\{−1,1\}$ be computed by a layered circuit of depth $d$, first-layer
width at most $w$, and tail size at most $s'$, where tail size means the number of
non-output gates in layers after the first, plus one. Then its total influence, computed
using the real-valued sign function, satisfies $\operatorname{Inf}(f)\le
2w\bigl(20(\lceil\log_2(s'+1)\rceil+2)\bigr)^{d-2}$.
\end{theorem}

\begin{theorem}[Total influence from circuit size]
\lean{FABL.DepthCircuit.exercise4_20b}
\proofsource{Analysis-of-Boolean-Functions-by-Ryan-ODonnell}{
  pdf-ffbe9c7b3306-p106-b007,
  pdf-ffbe9c7b3306-p110-b004,
  pdf-ffbe9c7b3306-p111-b003
}
\label{FABL.DepthCircuit.exercise4_20b}
\leanok
\uses{FABL.BooleanFunction, FABL.BooleanFunction.toReal, FABL.DepthCircuit.HasDepthSizeCircuit, FABL.DepthCircuit.exercise4_20b_of_quarterSwitching, FABL.DepthCircuit.initialRestrictionWidth, FABL.circuitInfluenceStep, FABL.hastadQuarterSwitchingBound, FABL.totalInfluence}
\difficulty{1}
Let $f:\{−1,1\}^n\to\{−1,1\}$ be computed by a layered circuit of depth $d$ and size at
most $s$, where size counts first-layer gates and all non-output gates in the remaining
layers. Then $\operatorname{Inf}(f)\le
4(5\lceil\log_2(s)\rceil+4)\bigl(20(\lceil\log_2(s+2)\rceil+2)\bigr)^{d-2}$.
\end{theorem}

\begin{theorem}[Influence bound for bounded-size circuits]
\lean{FABL.DepthCircuit.theorem4_30}
\proofsource{Analysis-of-Boolean-Functions-by-Ryan-ODonnell}{pdf-ffbe9c7b3306-p106-b007}
\label{FABL.DepthCircuit.theorem4_30}
\leanok
\uses{FABL.BooleanFunction, FABL.BooleanFunction.toReal, FABL.DepthCircuit.HasDepthSizeCircuit, FABL.DepthCircuit.theorem4_30_of_quarterSwitching, FABL.circuitInfluenceStep, FABL.hastadQuarterSwitchingBound, FABL.totalInfluence}
\difficulty{1}
Let $f$ be a Boolean function on the $n$-dimensional sign cube, and let $d,s$ be natural
numbers. If $f$ is computed by a depth-$d$ circuit of size at most $s$ for some width,
then its total influence is at most $4\bigl(20(\lceil\log_2(s+1)\rceil+2)\bigr)^{d-1}$.
\end{theorem}
